% Imporditud: tools/import_eksperimendid.py
% Allikas: Eesti füüsikaolümpiaadi eksperimendi ülesannete
%   kogu (Rael Kalda, 2023), ülesanne Ü29, lahendus L29.
\setAuthor{Oleg Košik}
\setRound{piirkonnavoor}
\setYear{2021}
\setNumber{P E1}
\setDifficulty{1}
\setTopic{vedelike mehaanika}
\setVahendid{tennisepall, joonlaud, liitrine anum veega ja õhukesest kartongist piima- või mahlapakk, mille ülemine tahk on ära lõigatud. Vee tiheduseks võib võtta 1000 kg m$^{-3}$}
\setPunktid{10}
\setAllikas{Eksperimendi kogu Ü29, lahendus lk 45}
\setLahendusseis{toores}

\prob{Tennisepall}
Leidke märja tennisepalli tihedus.

\hint

\solu
Tennisepalli massi saab määrata ujumise tingimusest

mpallg = $\rho$vgVvedelikualune = $\rho$vgS$\Delta$$h_1$,

kus S on piimapaki ristlõikepindala ja $\Delta$$h_1$ = $h_1$ $-$ $h_0$ veetaseme muutus. Tennisepalli ruumala määratakse sukeldusmeetodil surudes tennisepalli vee alla, millele vastab veetaseme muutus $\Delta$$h_2$ = $h_2$ $-$ $h_0$. Seega Vpall = S$\Delta$$h_2$. Kuna

$\rho$pall = mpall = $\rho$v S $\Delta$$h_1$ = $\rho$v $h_1$ $-$ $h_0$ , Vpall S$\Delta$$h_2$ $h_2$ $-$ $h_0$

taandub mõõtmine veetaseme kahe muutuse mõõtmisele.

Märkus. Tennisepalli läbimõõt on 6,70$\pm_0$,16 cm ning kuiva palli mass 6,70$\pm_0$,16 cm. Seetõttu kuiva palli tihedus jääb vahemikku 0,35 $-$ 0,39 g/cm$^3$, märja oma on suurusjärgus 0,45 $-$ 0,50 g/cm$^3$.

Hindamine

Kui on leitud valem veetaseme kahe muutuse jagatisena: Valemi tuletamine: 5 p. Mõõdetud $h_0$, $h_1$, $h_2$: 3 p. Arvutatud tihedus ja vastus realistlik: 2 p.

Kui on eraldi leitud mass ujumismeetodil ja ruumala sukeldusmeetodil: Mõõdetud aluse ristlõike küljepikkus ja arvutatud pindala: 1 p. Tuletatud valem ruumala jaoks sukeldusmeetodil: 1 p. Tehtud veetaseme muudu mõõtmised: 1 p. Arvutatud ruumala: 1 p. Tuletatud valem massi jaoks ujumismeetodil: 2 p. Tehtud veetaseme muudu mõõtmised: 1 p. Arvutatud mass: 1 p. Arvutatud tihedus ja vastus realistlik: 2 p.

Kui on eraldi leitud mass ujumismeetodil ja ruumala diameetri mõõtmisena (ebatäpne

meetod, kokku 8 p):

Ruumala valem V = 4 $\pi$$r_3$: 1 p. Hinnatud mõõtmise abil diameeter ja arvutatud ruumala: 1 p.

Mõõdetud aluse ristlõike küljepikkus ja arvutatud pindala: 1 p.

Tuletatud valem massi jaoks ujumismeetodil: 2 p.

Tehtud veetaseme muudu mõõtmised: 1 p.

Arvutatud mass: 1 p.

Arvutatud tihedus: 1 p.
\probend
